% =====================================================================
% Ch2 report — Convolutions + regularisation on Fashion MNIST (1 A4)
% =====================================================================
\documentclass[11pt,a4paper]{article}
\usepackage[margin=1.8cm]{geometry}
\usepackage{graphicx}
\usepackage{float}
\usepackage{caption}
\usepackage{booktabs}
\usepackage{fancyhdr}
\usepackage{parskip}
\setlength{\parskip}{2pt}
\usepackage{titlesec}
\titlespacing*{\section}{0pt}{3pt}{1pt}
\captionsetup{font=small,skip=2pt}
\setlength{\intextsep}{3pt}
\setlength{\textfloatsep}{3pt}

\newcommand{\studentname}{Vincent Brand}
\newcommand{\studentnumber}{1447806}
\newcommand{\course}{Machine Learning -- MADS, HAN}
\newcommand{\reportdate}{2026}
\newcommand{\repourl}{https://github.com/vincentbrand/MADS-ML-Vincent}

\pagestyle{fancy}
\fancyhf{}
\lhead{\studentname\ (\studentnumber)}
\rhead{\course}
\cfoot{\thepage}
\renewcommand{\headrulewidth}{0.4pt}
\sloppy

% include a figure if it exists, otherwise show a placeholder box.
% optional first arg = width (default \linewidth) so figures fit on 1 A4.
\newcommand{\figorbox}[2][\linewidth]{%
  \IfFileExists{#2}{\includegraphics[width=#1]{#2}}%
  {\fbox{\parbox[c][3.5cm][c]{0.98\linewidth}{\centering\ttfamily
   missing figure: #2\\[2pt] run the notebook to generate it}}}}

\begin{document}

\begin{center}
  {\Large\bfseries Convolutions, batch normalisation and dropout on Fashion MNIST}\\[3pt]
  \studentname\ \textbar\ \studentnumber\ \textbar\ \course\\
  \reportdate\ \textbar\ \repourl
\end{center}

\section{Hypothesis}
Ch1's dense network plateaued near $0.85$ --- extra width helped little, extra depth
not at all --- which motivates convolutions and regularisation. I test three
hypotheses. \textbf{(H1)} A small CNN should reach the dense baseline's accuracy at
\emph{far fewer} parameters, since convolution's parameter sharing and local
connectivity exploit the 2-D structure that flattening the 784 pixels destroys
(Ch10). \textbf{(H2)} BatchNorm should barely help a shallow CNN but stabilise a
deeper one and make it learning-rate tolerant, by smoothing a loss landscape whose
roughness only binds as depth grows (Ch11). \textbf{(H3)} Dropout should shrink the
train--validation gap and help accuracy \emph{only} when the model overfits; on an
under-capacity model it should \emph{hurt} by removing usable capacity (Ch9).

\section{Experiment design}
The CNN stacks conv blocks in an \texttt{nn.ModuleList} ---
\texttt{Conv2d}($3\times3$)$\to$ReLU$\to$optional
\texttt{BatchNorm2d}$\to$\texttt{MaxPool2d}, depth set by \texttt{num\_conv\_layers}
--- then an \texttt{AvgPool2d} and a dense head with optional \texttt{Dropout}; the
dense baseline is the Ch1 model. All runs share \texttt{CrossEntropyLoss}, Adam,
\texttt{ReduceLROnPlateau} and one fixed step budget for fair comparison, logged to
MLflow. H1 pits dense ($64\times32$ to $512\times256$) against CNNs (2--3 layers,
16--32 filters) by parameter count (Fig.~\ref{fig:h1}); H2 is a
depth$\times$BatchNorm$\times$lr grid (Fig.~\ref{fig:h2}); H3 varies dropout on a
small ($64\times32$) and a large ($512\times512$) model (Fig.~\ref{fig:h3}).

\section{Results}
\textbf{H1 only half held} (Fig.~\ref{fig:h1}): the CNNs are far more
parameter-efficient --- a 6k-param CNN already learns and the 24k-param 3-layer CNN
reaches $0.778$ at a $0.6\%$ gap --- yet none beat the dense nets on raw accuracy
(best dense $0.845$ at 536k params). The CNN curve still climbs steeply (the third
conv layer alone gave $0.678\!\to\!0.778$), so the binding limit is the narrow
32-filter extractor, not the principle. \textbf{H2 confirmed} (Fig.~\ref{fig:h2}):
BatchNorm helped in every cell, most where optimisation struggled (shallow,
lr$=10^{-3}$: $0.684\!\to\!0.838$); without it the net needed the higher lr, with it
lr$=10^{-3}$ gave the best model (deep+BN, $0.889$) --- BatchNorm \emph{decoupled}
accuracy from the learning rate. \textbf{H3 confirmed by regime}
(Fig.~\ref{fig:h3}): dropout $0.5$ shrank the large model's gap
($0.0159\!\to\!0.0133$) but on the small model left the gap flat and cut accuracy
($0.837\!\to\!0.783$).

\begin{figure}[H]
  \centering
  \figorbox[0.37\linewidth]{FigureH1.png}
  \caption{H1 --- validation accuracy vs.\ parameter count (log; up-and-left is
  better). CNNs use far fewer parameters but do not reach the dense nets' accuracy.}
  \label{fig:h1}
\end{figure}

\begin{figure}[H]
  \centering
  \begin{minipage}[t]{0.47\linewidth}
    \centering
    \figorbox[\linewidth]{FigureH2.png}
    \captionof{figure}{H2 --- valid.\ accuracy over conv depth $\times$ BatchNorm at
    two lr. BatchNorm helps everywhere and is nearly lr-independent.}
    \label{fig:h2}
  \end{minipage}\hfill
  \begin{minipage}[t]{0.47\linewidth}
    \centering
    \figorbox[\linewidth]{FigureH3.png}
    \captionof{figure}{H3 --- generalisation gap (left) and valid.\ accuracy (right),
    dropout $0.0$ vs.\ $0.5$. The gap shrinks only on the large model.}
    \label{fig:h3}
  \end{minipage}
\end{figure}

\section{Conclusions}
BatchNorm was the most valuable addition: it raised accuracy and made the learning
rate far less critical (the smoother-landscape effect of Ch11), with a payoff that
grew with depth. Convolutions gave efficiency and tighter generalisation but, at the
filter widths I could afford, did not overtake the dense baseline --- an honest
mismatch with H1 that points at the extractor, not the inductive bias. Dropout helped
only the overparameterised model, and only modestly, since Fashion-MNIST barely
overfits here. Next I would widen the CNN, drop the global-average head, add
augmentation, and test on data where regularisation bites harder.

\end{document}
